\documentclass[10pt,a4paper]{amsart}
\usepackage{amsmath,amssymb,amsthm,graphicx}
% For commutative diagrams we use Xy-pic
\usepackage[all]{xy}

% We use 2cell for 2-commutative diagrams.
\xyoption{2cell}
% Theorem environments.
%
\theoremstyle{plain}
\newtheorem{theorem}[subsection]{Theorem}
\newtheorem{proposition}[subsection]{Proposition}
\newtheorem{lemma}[subsection]{Lemma}

\theoremstyle{definition}
\newtheorem{definition}[subsection]{Definition}
\newtheorem{example}[subsection]{Example}
\newtheorem{exercise}[subsection]{Exercise}
\newtheorem{situation}[subsection]{Situation}

\theoremstyle{remark}
\newtheorem{remark}[subsection]{Remark}
\newtheorem{remarks}[subsection]{Remarks}

\numberwithin{equation}{subsection}

% Macros
%
\def\lim{\mathop{\mathrm{lim}}\nolimits}
\def\colim{\mathop{\mathrm{colim}}\nolimits}
\def\Spec{\mathop{\mathrm{Spec}}}
\def\Hom{\mathop{\mathrm{Hom}}\nolimits}
\def\Ext{\mathop{\mathrm{Ext}}\nolimits}
\def\SheafHom{\mathop{\mathcal{H}\!\mathit{om}}\nolimits}
\def\SheafExt{\mathop{\mathcal{E}\!\mathit{xt}}\nolimits}
\def\Sch{\mathit{Sch}}
\def\Mor{\mathop{\mathrm{Mor}}\nolimits}
\def\Ob{\mathop{\mathrm{Ob}}\nolimits}
\def\Sh{\mathop{\mathit{Sh}}\nolimits}
\def\NL{\mathop{N\!L}\nolimits}
\def\CH{\mathop{\mathrm{CH}}\nolimits}
\def\proetale{{pro\text{-}\acute{e}tale}}
\def\etale{{\acute{e}tale}}
\def\QCoh{\mathit{QCoh}}
\def\Ker{\mathop{\mathrm{Ker}}}
\def\Im{\mathop{\mathrm{Im}}}
\def\Coker{\mathop{\mathrm{Coker}}}
\def\Coim{\mathop{\mathrm{Coim}}}

% Boxtimes
%
\DeclareMathSymbol{\boxtimes}{\mathbin}{AMSa}{"02}

%
% Macros for moduli stacks/spaces
%
\def\QCohstack{\mathcal{QC}\!\mathit{oh}}
\def\Cohstack{\mathcal{C}\!\mathit{oh}}
\def\Spacesstack{\mathcal{S}\!\mathit{paces}}
\def\Quotfunctor{\mathrm{Quot}}
\def\Hilbfunctor{\mathrm{Hilb}}
\def\Curvesstack{\mathcal{C}\!\mathit{urves}}
\def\Polarizedstack{\mathcal{P}\!\mathit{olarized}}
\def\Complexesstack{\mathcal{C}\!\mathit{omplexes}}
% \Pic is the operator that assigns to X its picard group, usage \Pic(X)
% \Picardstack_{X/B} denotes the Picard stack of X over B
% \Picardfunctor_{X/B} denotes the Picard functor of X over B
\def\Pic{\mathop{\mathrm{Pic}}\nolimits}
\def\Picardstack{\mathcal{P}\!\mathit{ic}}
\def\Picardfunctor{\mathrm{Pic}}
\def\Deformationcategory{\mathcal{D}\!\mathit{ef}}
\usepackage{cleveref}
\crefname{originallemma}{original lemma}{original lemmas}
\Crefname{originallemma}{Original lemma}{Original lemmas}
\makeatletter
\newlabel{topologies-definition-standard-Zariski}{{020R}{source}{}{sourcepointer-1}{}}
\newlabel{topologies-definition-standard-Zariski@cref}{{[originallemma][0][]020R}{[1][0][]source}{}{}{}}
\makeatother
\crefname{originallemma}{original lemma}{original lemmas}
\Crefname{originallemma}{Original lemma}{Original lemmas}
\makeatletter
\newlabel{topologies-definition-standard-etale}{{0219}{source}{}{sourcepointer-2}{}}
\newlabel{topologies-definition-standard-etale@cref}{{[originallemma][0][]0219}{[1][0][]source}{}{}{}}
\makeatother
\crefname{originallemma}{original lemma}{original lemmas}
\Crefname{originallemma}{Original lemma}{Original lemmas}
\makeatletter
\newlabel{topologies-definition-standard-smooth}{{0223}{source}{}{sourcepointer-3}{}}
\newlabel{topologies-definition-standard-smooth@cref}{{[originallemma][0][]0223}{[1][0][]source}{}{}{}}
\makeatother
\crefname{originallemma}{original lemma}{original lemmas}
\Crefname{originallemma}{Original lemma}{Original lemmas}
\makeatletter
\newlabel{topologies-definition-standard-syntomic}{{0229}{source}{}{sourcepointer-4}{}}
\newlabel{topologies-definition-standard-syntomic@cref}{{[originallemma][0][]0229}{[1][0][]source}{}{}{}}
\makeatother
\crefname{originallemma}{original lemma}{original lemmas}
\Crefname{originallemma}{Original lemma}{Original lemmas}
\makeatletter
\newlabel{topologies-definition-standard-fppf}{{021Q}{source}{}{sourcepointer-5}{}}
\newlabel{topologies-definition-standard-fppf@cref}{{[originallemma][0][]021Q}{[1][0][]source}{}{}{}}
\makeatother
\crefname{originalremark}{original remark}{original remarks}
\Crefname{originalremark}{Original remark}{Original remarks}
\makeatletter
\newlabel{example-fpqc-coverings}{{03NY}{source}{}{sourcepointer-6}{}}
\newlabel{example-fpqc-coverings@cref}{{[originalremark][0][]03NY}{[1][0][]source}{}{}{}}
\makeatother
\crefname{originallemma}{original lemma}{original lemmas}
\Crefname{originallemma}{Original lemma}{Original lemmas}
\makeatletter
\newlabel{topologies-definition-etale-covering}{{0215}{source}{}{sourcepointer-7}{}}
\newlabel{topologies-definition-etale-covering@cref}{{[originallemma][0][]0215}{[1][0][]source}{}{}{}}
\makeatother
\crefname{originallemma}{original lemma}{original lemmas}
\Crefname{originallemma}{Original lemma}{Original lemmas}
\makeatletter
\newlabel{topologies-definition-smooth-covering}{{021Z}{source}{}{sourcepointer-8}{}}
\newlabel{topologies-definition-smooth-covering@cref}{{[originallemma][0][]021Z}{[1][0][]source}{}{}{}}
\makeatother
\crefname{originallemma}{original lemma}{original lemmas}
\Crefname{originallemma}{Original lemma}{Original lemmas}
\makeatletter
\newlabel{topologies-definition-syntomic-covering}{{0225}{source}{}{sourcepointer-9}{}}
\newlabel{topologies-definition-syntomic-covering@cref}{{[originallemma][0][]0225}{[1][0][]source}{}{}{}}
\makeatother
\crefname{originallemma}{original lemma}{original lemmas}
\Crefname{originallemma}{Original lemma}{Original lemmas}
\makeatletter
\newlabel{topologies-definition-fppf-covering}{{021M}{source}{}{sourcepointer-10}{}}
\newlabel{topologies-definition-fppf-covering@cref}{{[originallemma][0][]021M}{[1][0][]source}{}{}{}}
\makeatother
\crefname{originallemma}{original lemma}{original lemmas}
\Crefname{originallemma}{Original lemma}{Original lemmas}
\makeatletter
\newlabel{sites-cohomology-equation-map-cech-complexes}{{03F4}{source}{}{sourcepointer-11}{}}
\newlabel{sites-cohomology-equation-map-cech-complexes@cref}{{[originallemma][0][]03F4}{[1][0][]source}{}{}{}}
\makeatother
\title{A standard covering is refined by a single fpqc covering with the same Cech complex}
\author{The Stacks Project Authors}
\date{stacks-project:etale-cohomology.tex:local-snapshot-20261011; GNU FDL 1.2 or later}
\begin{document}
\maketitle
\noindent\textit{Editorial note.} All delivered mathematical text is the author's own, copied from the pinned Stacks Project chapter file etale-cohomology.tex; no statement or proof step was rewritten, reordered or removed.\par\medskip
\noindent\textit{Editorial note.} The proof is the author's complete argument for the three assertions; the refinement maps in both directions are shown to be inverse, which settles the third.\par\medskip
\noindent\textit{Editorial note.} Every labelled reference occurring inside the counted spans is carried in assembly.external\_references as a non-load-bearing navigation pointer (same-chapter ones at their definition site, cross-chapter ones at the invocation inside the counted span), so no load-bearing mathematics is replaced by a stub. No bibliography citation occurs inside the counted spans.\par\medskip
\noindent\textit{Editorial note.} The checked chapter list for same-name collisions is the set of stacks chapters staged in this repository: chow, cohomology, divisors, etale-cohomology, flat, more-algebra, more-morphisms, morphisms, spaces-more-morphisms, varieties together with this card's chapter etale-cohomology.\par\medskip
\noindent\textit{Editorial note.} The preamble is the chapter's own native preamble excerpt already used for this source: the Xy-pic setup (native preamble lines 14-18) and the native theorem declarations with the mathematical macros (native preamble lines 154-217).\par\medskip
\noindent\textit{Allowed standard prerequisites (editorial).}\begin{itemize}
\item Standard tau-covering and the big or small site of a scheme. Same chapter, the definition of a standard tau-covering of the Zariski, etale, smooth, syntomic or fppf site of a scheme. Supplies the covering U and the sites in which the statement is made.
\item Cech complex of a covering and comparison of Cech complexes under refinement. Same chapter, the Cech complex of a covering and the map induced by a refinement. Supplies the third assertion of the statement and the strategy of comparing the two complexes through refinement maps.
\end{itemize}\medskip
\ifdefined\hypertarget\hypertarget{sourcepointer-1}{}\fi
\noindent\textit{Source reference.} Original lemma 020R: sources/stacks-etale-cohomology/pinned-original-source.tex, source lines 2188--2188. Navigation pointer invoked at proof line 2188; the counted span contains the pointer only and the cited result is not reproduced inside it, so the pointer is registered as a non-load-bearing navigation stub in the established form for Stacks material. No bibliography citation is involved.\par\medskip
\ifdefined\hypertarget\hypertarget{sourcepointer-2}{}\fi
\noindent\textit{Source reference.} Original lemma 0219: sources/stacks-etale-cohomology/pinned-original-source.tex, source lines 2189--2189. Navigation pointer invoked at proof line 2189; the counted span contains the pointer only and the cited result is not reproduced inside it, so the pointer is registered as a non-load-bearing navigation stub in the established form for Stacks material. No bibliography citation is involved.\par\medskip
\ifdefined\hypertarget\hypertarget{sourcepointer-3}{}\fi
\noindent\textit{Source reference.} Original lemma 0223: sources/stacks-etale-cohomology/pinned-original-source.tex, source lines 2190--2190. Navigation pointer invoked at proof line 2190; the counted span contains the pointer only and the cited result is not reproduced inside it, so the pointer is registered as a non-load-bearing navigation stub in the established form for Stacks material. No bibliography citation is involved.\par\medskip
\ifdefined\hypertarget\hypertarget{sourcepointer-4}{}\fi
\noindent\textit{Source reference.} Original lemma 0229: sources/stacks-etale-cohomology/pinned-original-source.tex, source lines 2191--2191. Navigation pointer invoked at proof line 2191; the counted span contains the pointer only and the cited result is not reproduced inside it, so the pointer is registered as a non-load-bearing navigation stub in the established form for Stacks material. No bibliography citation is involved.\par\medskip
\ifdefined\hypertarget\hypertarget{sourcepointer-5}{}\fi
\noindent\textit{Source reference.} Original lemma 021Q: sources/stacks-etale-cohomology/pinned-original-source.tex, source lines 2192--2192. Navigation pointer invoked at proof line 2192; the counted span contains the pointer only and the cited result is not reproduced inside it, so the pointer is registered as a non-load-bearing navigation stub in the established form for Stacks material. No bibliography citation is involved.\par\medskip
\ifdefined\hypertarget\hypertarget{sourcepointer-6}{}\fi
\noindent\textit{Source reference.} Original remark 03NY: sources/stacks-etale-cohomology/pinned-original-source.tex, source lines 905--905. Navigation pointer invoked at proof line 2197; the counted span contains the pointer only and the cited result is not reproduced inside it, so the pointer is registered as a non-load-bearing navigation stub in the established form for Stacks material. No bibliography citation is involved.\par\medskip
\ifdefined\hypertarget\hypertarget{sourcepointer-7}{}\fi
\noindent\textit{Source reference.} Original lemma 0215: sources/stacks-etale-cohomology/pinned-original-source.tex, source lines 2201--2201. Navigation pointer invoked at proof line 2201; the counted span contains the pointer only and the cited result is not reproduced inside it, so the pointer is registered as a non-load-bearing navigation stub in the established form for Stacks material. No bibliography citation is involved.\par\medskip
\ifdefined\hypertarget\hypertarget{sourcepointer-8}{}\fi
\noindent\textit{Source reference.} Original lemma 021Z: sources/stacks-etale-cohomology/pinned-original-source.tex, source lines 2202--2202. Navigation pointer invoked at proof line 2202; the counted span contains the pointer only and the cited result is not reproduced inside it, so the pointer is registered as a non-load-bearing navigation stub in the established form for Stacks material. No bibliography citation is involved.\par\medskip
\ifdefined\hypertarget\hypertarget{sourcepointer-9}{}\fi
\noindent\textit{Source reference.} Original lemma 0225: sources/stacks-etale-cohomology/pinned-original-source.tex, source lines 2203--2203. Navigation pointer invoked at proof line 2203; the counted span contains the pointer only and the cited result is not reproduced inside it, so the pointer is registered as a non-load-bearing navigation stub in the established form for Stacks material. No bibliography citation is involved.\par\medskip
\ifdefined\hypertarget\hypertarget{sourcepointer-10}{}\fi
\noindent\textit{Source reference.} Original lemma 021M: sources/stacks-etale-cohomology/pinned-original-source.tex, source lines 2204--2204. Navigation pointer invoked at proof line 2204; the counted span contains the pointer only and the cited result is not reproduced inside it, so the pointer is registered as a non-load-bearing navigation stub in the established form for Stacks material. No bibliography citation is involved.\par\medskip
\ifdefined\hypertarget\hypertarget{sourcepointer-11}{}\fi
\noindent\textit{Source reference.} Original lemma 03F4: sources/stacks-etale-cohomology/pinned-original-source.tex, source lines 2212--2212. Navigation pointer invoked at proof line 2212; the counted span contains the pointer only and the cited result is not reproduced inside it, so the pointer is registered as a non-load-bearing navigation stub in the established form for Stacks material. No bibliography citation is involved.\par\medskip
\begin{lemma}
\label{lemma-cech-complex}
Let $\tau \in \{fppf, syntomic, smooth, \etale, Zariski\}$.
Let $S$ be a scheme.
Let $\mathcal{F}$ be an abelian sheaf on $(\Sch/S)_\tau$, or on
$S_\tau$ in case $\tau = \etale$, and let
$\mathcal{U} = \{U_i \to U\}_{i \in I}$
be a standard $\tau$-covering of this site.
Let $V = \coprod_{i \in I} U_i$. Then
\begin{enumerate}
\item $V$ is an affine scheme,
\item $\mathcal{V} = \{V \to U\}$ is an fpqc covering
and also a $\tau$-covering unless $\tau = Zariski$,
\item the {\v C}ech complexes
$\check{\mathcal{C}}^\bullet (\mathcal{U}, \mathcal{F})$ and
$\check{\mathcal{C}}^\bullet (\mathcal{V}, \mathcal{F})$ agree.
\end{enumerate}
\end{lemma}

\begin{proof}
The definition of a standard $\tau$-covering is given in
Topologies, Definition
\ref{topologies-definition-standard-Zariski},
\ref{topologies-definition-standard-etale},
\ref{topologies-definition-standard-smooth},
\ref{topologies-definition-standard-syntomic}, and
\ref{topologies-definition-standard-fppf}.
By definition each of the schemes
$U_i$ is affine and $I$ is a finite set. Hence $V$ is an affine scheme.
It is clear that $V \to U$ is flat and surjective, hence
$\mathcal{V}$ is an fpqc covering, see
Example \ref{example-fpqc-coverings}.
Excepting the Zariski case, the covering $\mathcal{V}$
is also a $\tau$-covering, see
Topologies, Definition
\ref{topologies-definition-etale-covering},
\ref{topologies-definition-smooth-covering},
\ref{topologies-definition-syntomic-covering}, and
\ref{topologies-definition-fppf-covering}.

\medskip\noindent
Note that $\mathcal{U}$ is a refinement of $\mathcal{V}$
and hence there is a map of {\v C}ech complexes
$\check{\mathcal{C}}^\bullet (\mathcal{V}, \mathcal{F}) \to
\check{\mathcal{C}}^\bullet (\mathcal{U}, \mathcal{F})$, see
Cohomology on Sites,
Equation (\ref{sites-cohomology-equation-map-cech-complexes}).
Next, we observe that if $T = \coprod_{j \in J} T_j$ is a
disjoint union of schemes in the site on which $\mathcal{F}$ is defined
then the family of morphisms with fixed target
$\{T_j \to T\}_{j \in J}$ is a Zariski covering, and so
\begin{equation}
\label{equation-sheaf-coprod}
\mathcal{F}(T) =
\mathcal{F}(\coprod\nolimits_{j \in J} T_j) =
\prod\nolimits_{j \in J} \mathcal{F}(T_j)
\end{equation}
by the sheaf condition of $\mathcal{F}$.
This implies the map of {\v C}ech complexes above is an isomorphism
in each degree because
$$
V \times_U \ldots \times_U V
=
\coprod\nolimits_{i_0, \ldots i_p} U_{i_0} \times_U \ldots \times_U U_{i_p}
$$
as schemes.
\end{proof}

\end{document}
